\chapter{ETFs, Futures and Index Arbitrage}\label{ch:iv2:etfs-futures-and-index-arbitrage}

\begin{ivtrains}
The delta-one desk's arithmetic: what an ETF should trade at, when creating or redeeming pays, what a future's price contains, what a leveraged fund's daily reset does to it and to the market, and how much a rebalance or a roll is worth. Rates here are simple unless a stem says otherwise.
\end{ivtrains}

\begin{question}[$\star$ \iqroles{QT, Trader} \iqfirm{market maker, hedge fund, bank} \iqlevel{junior}]\label{iq:iv2:etfs-futures-and-index-arbitrage:1}
For every 100 of its shares, an ETF holds 50 shares of a stock at \$20, 30 of a stock at \$50, 10 of a stock at \$100, and \$500 in cash. What is its net asset value per share? It trades at \$40.10: what is its premium in basis points?
\end{question}

\begin{question}[$\star$ \iqroles{QT, Trader} \iqfirm{market maker, hedge fund} \iqlevel{junior}]\label{iq:iv2:etfs-futures-and-index-arbitrage:2}
An ETF's net asset value is \$40.00. Buying the basket and creating costs a total of 6 cents per ETF share, and redeeming and selling the basket costs a total of 4 cents. Between what ETF prices is there no arbitrage? The bid is \$40.09: what do you do, and how much do you make per share?
\end{question}

\begin{question}[$\star$ \iqroles{QT, Trader, QR} \iqfirm{bank, hedge fund, market maker} \iqlevel{junior}]\label{iq:iv2:etfs-futures-and-index-arbitrage:3}
A stock trades at \$100. It pays a \$1 dividend in three months. Interest is 4\% a year, and you borrow and lend at that rate. What is the fair price of a six-month forward? Show the trades that lock in a profit if the forward trades at \$102.
\end{question}

\begin{question}[$\star$ \iqroles{QT, Trader, Risk} \iqfirm{hedge fund, bank} \iqlevel{junior}]\label{iq:iv2:etfs-futures-and-index-arbitrage:4}
An index stands at 4\,000, and its future has a multiplier of \$50 per point. What is one contract's notional value? How many contracts hedge a \$10 million portfolio with a beta of 1, and how many with a beta of 1.2?
\end{question}

\begin{question}[$\star$ \iqroles{QT, Trader} \iqfirm{bank, hedge fund, market maker} \iqlevel{junior}]\label{iq:iv2:etfs-futures-and-index-arbitrage:5}
An index stands at 4\,000. Its future, expiring in a quarter of a year, trades at 4\,020. The dividend yield is 1\% a year. What financing rate does the future imply? If you can borrow at 2.5\%, what would you do?
\end{question}

\begin{question}[$\star$ \iqroles{QT, Trader} \iqfirm{bank, hedge fund, market maker} \iqlevel{junior}]\label{iq:iv2:etfs-futures-and-index-arbitrage:6}
An index stands at 4\,000, and its quarter-year future trades at 4\,030. The interest rate is 5\% a year. What dividends, in index points and as an annual yield, does the future imply?
\end{question}

\begin{question}[$\star$ \iqroles{QT, Trader, QR, Analyst} \iqfirm{any} \iqlevel{junior}]\label{iq:iv2:etfs-futures-and-index-arbitrage:7}
An index rises 2\% and then falls 2\%. What is its total return over the two days? A fund resets each day to three times the index's daily return. What is its return? How does that compare with three times the index's return?
\end{question}

\begin{question}[$\star$ \iqroles{QT, Trader, Analyst} \iqfirm{any} \iqlevel{junior}]\label{iq:iv2:etfs-futures-and-index-arbitrage:8}
An index goes from 100 to 110 and back to 100. A fund resets each day to minus one times the index's daily return. What is the fund's total return? Someone holds it to hedge a long index position for those two days. What is the combined result?
\end{question}

\begin{question}[$\star$ \iqroles{QT, QR, Analyst} \iqfirm{hedge fund, bank} \iqlevel{junior}]\label{iq:iv2:etfs-futures-and-index-arbitrage:9}
A fund holds two stocks equally, \$100 million in each. One rises 10\%, and the other falls 10\%. What trades bring it back to equal weights, and what is the one-way turnover as a share of the fund? What trades would a cap-weighted fund need? Why?
\end{question}

\begin{question}[$\star\star$ \iqroles{QT, Trader, Risk} \iqfirm{hedge fund, bank} \iqlevel{junior}]\label{iq:iv2:etfs-futures-and-index-arbitrage:10}
You are long 10 index futures at 4\,000 with a \$50 multiplier, and you post 5\% of notional as initial margin. The index falls to 3\,900 by the close. What is your loss, and how much cash moves today? What is your margin if the initial margin is still required at the new price? What is the difference between this position and the same exposure in a forward?
\end{question}

\begin{question}[$\star\star$ \iqroles{QT, Trader, Analyst} \iqfirm{any} \iqlevel{junior}]\label{iq:iv2:etfs-futures-and-index-arbitrage:11}
A fund stays long by holding the nearest monthly oil future and rolling it each month. The front future trades at 70 and the next at 71.40. Spot and the whole curve stay exactly the same for a year. What does the fund earn from the roll in a month? What over a year? Why does an unchanged oil price not give it an unchanged result?
\end{question}

\begin{question}[$\star\star$ \iqroles{QT, Trader, QR} \iqfirm{market maker, hedge fund} \iqlevel{junior, experienced}]\label{iq:iv2:etfs-futures-and-index-arbitrage:12}
A fund with \$1 billion of assets resets each day to three times the index. The index rises 1\% today. What are its assets and exposure before the reset, and how much must it buy or sell at the close? Write a formula for a leverage $L$, a fund size $A$ and a return $r$. What does a minus-three fund do on the same day?
\end{question}

\begin{question}[$\star\star$ \iqroles{QR, QT} \iqfirm{hedge fund, market maker} \iqlevel{junior, experienced}]\label{iq:iv2:etfs-futures-and-index-arbitrage:13}
A fund resets daily to $L$ times an index with annual volatility $\sigma$, with no rates or fees. Over a year, how does its return relate to the index's? If the index has 30\% volatility and finishes the year flat, what do the three-times fund and the minus-one fund return?
\end{question}

\begin{question}[$\star\star$ \iqroles{QT, QR, Dev} \iqfirm{bank, hedge fund} \iqlevel{junior}]\label{iq:iv2:etfs-futures-and-index-arbitrage:14}
A price-weighted index has three stocks at \$100, \$50 and \$30, and a divisor of 3, so it stands at 60. The first stock splits 2 for 1. What must the new divisor be, and why? If all three stocks then rise 10\%, what is the index, and what would it have been without the split?
\end{question}

\begin{question}[$\star\star$ \iqroles{QT, Trader} \iqfirm{bank, hedge fund, market maker} \iqlevel{junior}]\label{iq:iv2:etfs-futures-and-index-arbitrage:15}
A price index stands at 4\,000. A stock with a 2\% weight goes ex-dividend today, paying 1\% of its price. How many points does the index fall at the open, with all else unchanged? What happens to a total-return version of the index, and to the fair value of the future expiring after the date?
\end{question}

\begin{question}[$\star\star$ \iqroles{QT, Trader} \iqfirm{bank, hedge fund} \iqlevel{junior, experienced}]\label{iq:iv2:etfs-futures-and-index-arbitrage:16}
An index stands at 4\,000. You buy the basket and sell a one-year future at a price assuming dividends of 2\% of spot. Dividends turn out to be 2.5\%. What is your gain in points? Who is on the other side of this risk? What would you trade if you only wanted this dividend exposure?
\end{question}

\begin{question}[$\star\star$ \iqroles{QT, Trader, Analyst} \iqfirm{hedge fund, bank} \iqlevel{junior, experienced}]\label{iq:iv2:etfs-futures-and-index-arbitrage:17}
You want index exposure. A rolling future costs 15 basis points a year of implied financing above your funding rate, plus 4 basis points a year in roll costs. An ETF has a 9-basis-point annual fee, earns 2 basis points a year from lending, and costs 1 basis point in spread to buy and sell. Over one year, which is cheaper? At what holding period are they equal?
\end{question}

\begin{question}[$\star\star$ \iqroles{QT, Trader} \iqfirm{bank, hedge fund, market maker} \iqlevel{junior, experienced}]\label{iq:iv2:etfs-futures-and-index-arbitrage:18}
An index future expiring in three months trades at 4\,020, and the next one, three months later, trades at 4\,050. Rates are 5\% a year, and dividends over that later quarter run at 2\% a year. Is the calendar spread fair, using 4\,020 as the base? What if the later future trades at 4\,060? What do you do if you are long the near one and must stay long?
\end{question}

\begin{question}[$\star\star$ \iqroles{Analyst, QT} \iqfirm{hedge fund, bank} \iqlevel{junior}]\label{iq:iv2:etfs-futures-and-index-arbitrage:19}
An index fund has a fee of 0.20\% a year and earns 0.08\% from lending. It keeps 1\% of its assets in cash, and over the year the index returns 8\% while cash returns 4\%. It holds a sample of the index whose return misses it by 5 basis points a year, with mean zero. What is its expected return gap with the index? Distinguish tracking difference from tracking error.
\end{question}

\begin{question}[$\star\star$ \iqroles{QT, Trader} \iqfirm{bank, hedge fund} \iqlevel{junior, experienced}]\label{iq:iv2:etfs-futures-and-index-arbitrage:20}
An index stands at 4\,000, with a half-year future and dividends of 2\% a year. You borrow at 5\% and lend at 4.5\%. Selling the basket short costs a further fee of 0.3\% a year. Entering and leaving the trade costs 0.1\% of spot in total. Between what future prices is there no arbitrage for you? Why is the band not symmetric?
\end{question}

\begin{question}[$\star\star$ \iqroles{QT, Trader, Risk} \iqfirm{market maker, hedge fund, bank} \iqlevel{junior, experienced}]\label{iq:iv2:etfs-futures-and-index-arbitrage:21}
In a credit sell-off, a bond ETF trades at 95 while its published net asset value is 100. You would receive bonds that, sold right now, would really bring 95.50, and it costs 0.20 to redeem. What do you make per ETF share if everything happens at once? Selling the bonds takes three days, with a daily volatility of 0.5\% in their price. Is it arbitrage? Which number is the better estimate of value, the 95 or the 100?
\end{question}

\begin{question}[$\star\star$ \iqroles{QT, Trader} \iqfirm{market maker} \iqlevel{junior, experienced}]\label{iq:iv2:etfs-futures-and-index-arbitrage:22}
You quote an ETF at \$40.00 and \$40.02, fair value being \$40.01. When filled, you hedge at once with a future, paying half of its spread, 0.5 basis points. If your hedge and your ETF move apart while you hold them, what is the risk? What do you net per share on one fill? For 100\,000 shares what is that in dollars, and what does a basis move of 1 basis point against you do?
\end{question}

\begin{question}[$\star\star$ \iqroles{QT, Trader, Analyst} \iqfirm{market maker, hedge fund} \iqlevel{junior, experienced}]\label{iq:iv2:etfs-futures-and-index-arbitrage:23}
Buying an ETF's basket costs half a spread of 2 basis points. Over a day, investors buy \$100 million of the ETF and sell \$90 million, and a market maker takes the other side of all of it. What must the market maker hedge in the basket? What does that cost, per dollar of ETF traded? Why can an ETF's spread be narrower than the spread of what it holds?
\end{question}

\begin{question}[$\star\star$ \iqroles{QT, Trader, Analyst} \iqfirm{bank, hedge fund} \iqlevel{junior, experienced}]\label{iq:iv2:etfs-futures-and-index-arbitrage:24}
You want exposure to a stock yielding 2\% a year. Holding it in cash costs you funding plus 30 basis points, and you lose 15\% of the dividend to withholding tax. Holding it through a total return swap costs funding plus 40 basis points, and the dealer passes on 95\% of the dividend. Which is cheaper, and by how much? What risk do you take on with the swap, and what does the dealer do to hedge?
\end{question}

\begin{question}[$\star\star\star$ \iqroles{QT, Trader, Analyst} \iqfirm{hedge fund, market maker} \iqlevel{junior, experienced}]\label{iq:iv2:etfs-futures-and-index-arbitrage:25}
A stock with a market value of \$10 billion, 80\% of it free float, joins a cap-weighted index. Funds that track the index together own 20\% of the free float of every member. What must they buy? Daily volume is \$100 million: how many days is that at 100\% of volume, and at 20\%? With the rule that the cost is $\sigma\sqrt{Q/V}$ and 2\% daily volatility, what does the price do, and why should you doubt that number?
\end{question}

\begin{question}[$\star\star\star$ \iqroles{QR, QT} \iqfirm{hedge fund, market maker} \iqlevel{junior, experienced}]\label{iq:iv2:etfs-futures-and-index-arbitrage:26}
A colleague notices that a three-times and a minus-three-times fund both lose value in volatile markets, and proposes selling \$1 of each short and holding for a year, ignoring borrow costs. Daily returns have mean zero and an annual volatility of 30\%. What is your expected profit? What is your median profit, and how often do you lose money? What does this say about the ``decay''?
\end{question}

\begin{question}[$\star\star\star$ \iqroles{QT, QR, Trader} \iqfirm{market maker, hedge fund} \iqlevel{junior, experienced}]\label{iq:iv2:etfs-futures-and-index-arbitrage:27}
Three-times funds hold \$10 billion on an index, and minus-three-times funds hold \$5 billion. The index falls 4\% today. How much does each group have to trade at the close, and in which direction? Check each by computing exposures. If the market's liquidity in the last 15 minutes is \$20 billion, what share is this? Why does this look like the hedge of a short gamma position?
\end{question}

\begin{question}[$\star\star\star$ \iqroles{QR, QT} \iqfirm{bank, hedge fund} \iqlevel{experienced}]\label{iq:iv2:etfs-futures-and-index-arbitrage:28}
A stock pays nothing and is at 100. The one-period rate is 4\%. Under the risk-neutral measure, next period it is either 120 with a one-period rate of 6\%, or 84 with a one-period rate of 2\%, and in each state it then grows at that rate for one more period. Find the risk-neutral probabilities. What are the two-period forward price and the two-period futures price? Why are they different, and what makes the difference zero?
\end{question}

\begin{question}[$\star\star\star$ \iqroles{QR, QT, Risk} \iqfirm{hedge fund, bank} \iqlevel{junior, experienced}]\label{iq:iv2:etfs-futures-and-index-arbitrage:29}
A sector ETF has 18\% annual volatility. You hedge it with two stocks, $A$ of 20\% and $B$ of 30\% volatility, with a correlation of 0.5 with each other. The ETF has correlations of 0.8 with $A$ and 0.7 with $B$. What minimum-variance hedge weights do you use per dollar of ETF, and what is the remaining volatility? Why not hedge with 100 names?
\end{question}

\begin{question}[$\star\star\star$ \iqroles{QR, QT, Trader} \iqfirm{market maker, hedge fund} \iqlevel{junior, experienced}]\label{iq:iv2:etfs-futures-and-index-arbitrage:30}
A future is 6 points rich to fair value, and it must converge to fair value at expiry, in three days. The gap moves up or down by 4 points each day with equal probability on days 1 and 2, then goes to 0 on day 3. You may hold a position of $+1$, $0$ or $-1$ in the gap. If you stay short from today until expiry, what do you make? What is the most you can expect with free flipping? What is that option worth if each flip costs 0.5 points in total? When do you use it?
\end{question}

\begin{question}[$\star\star\star$ \iqroles{QT, Trader, Risk} \iqfirm{hedge fund, market maker} \iqlevel{junior, experienced}]\label{iq:iv2:etfs-futures-and-index-arbitrage:31}
A commodity fund stops issuing new shares because it has reached a limit on its holdings, while the old shares keep trading. Demand is strong, and the shares reach a 20\% premium to net asset value. You buy at that level. A month later the underlying is up 5\%, and creation resumes, so the premium disappears. What do you return? Why does the normal arbitrage not close the premium? What would you do with a short?
\end{question}

\begin{question}[$\star\star\star$ \iqroles{QT, Trader, Risk, Analyst} \iqfirm{hedge fund, bank} \iqlevel{junior, experienced}]\label{iq:iv2:etfs-futures-and-index-arbitrage:32}
An ETF has 100 million shares outstanding and a short interest of 120 million shares. How can more than 100\% be sold short? How many shares do holders of record and those long through loans believe they own in total? What do market makers do when they deliver? What is the actual risk to investors if the ETF closes?
\end{question}

\begin{question}[$\star\star\star$ \iqroles{Analyst, QT} \iqfirm{hedge fund, bank} \iqlevel{junior, experienced}]\label{iq:iv2:etfs-futures-and-index-arbitrage:33}
A fund tracks a cap-weighted index. A stock joins it with a weight of 0.1\%. Buying at the effective close is expected to cost 3\% in impact on that stock. Buying over the five days before costs 1\%, but the fund then holds an extra position, which you should take as the whole 0.1\% for all five days. The stock has a daily volatility of 2\%, hedged only by the rest of the index. What is the expected saving for the whole fund, in basis points? What is the tracking risk it takes on? Which would you do?
\end{question}

\begin{question}[$\star\star\star\star$ \iqroles{QR} \iqfirm{hedge fund, market maker, bank} \iqlevel{experienced}]\label{iq:iv2:etfs-futures-and-index-arbitrage:34}
A fund targets $L$ times the return on an index with constant volatility $\sigma$ and rate $r$, continuously rebalanced, with fee $f$. Using Ito's lemma, derive its value at time $T$ in terms of the index. For $L = 3$, $\sigma = 20\%$, $T = 1$, no fee and a fund starting at 1, what is its value if the index is unchanged and rates are zero, and what if the index rises 10\% and rates are 4\%? With zero rates, what is an at-the-money one-year call on the fund worth, per 100 of fund value? Compare it with three one-year at-the-money calls on an index at 100.
\end{question}

\begin{question}[$\star\star\star\star$ \iqroles{QT, QR, Trader} \iqfirm{hedge fund} \iqlevel{experienced}]\label{iq:iv2:etfs-futures-and-index-arbitrage:35}
Funds must buy \$1.6 billion of a stock at the close in five trading days, on an announced addition to an index. Daily volume is \$100 million and daily volatility is 2\%. Use a peak price pressure of $\sigma\sqrt{Q/V}$. Half of it was priced on the announcement. You buy \$100 million over the five days at 20\% of volume, at an average cost of $\sigma\sqrt{Q/V}$ on your own size, and you sell to the funds at the close. What is your expected profit, and what risk do you run? If others own \$400 million by the close and do the same, what happens to the pressure? When would you pass?
\end{question}
